\documentclass{ximera}
\title{Solar Days}

\newcommand{\pskip}{\vskip 0.1 in}

\begin{document}
\begin{abstract}
The length of a solar day.
\end{abstract}
\maketitle



Most days do not have 24 hours.

Twenty four hours/day is the length of the \emph{mean} (ie. average) \emph{solar day}, and the length of the solar day varies about this mean through the course of a year.

A solar day is the time between consecutive \emph{solar noons}.  Solar noon occurs when the sun crosses the local meridian and is at its highest point in the sky (due south here in Shoreline). As we approach the winter solstice, the solar days are approaching their maximum length, and exceed 24 hours by as much as 30 seconds. The solar days are the shortest around March 26, when they fall about 18 seconds short of the 24-hour mean.

There are two reasons for the variation of the length of the solar day about its mean. One is the \emph{obliquity} of the earth's orbit about the sun. The obliquity, also called \emph{axial tilt} and currently about $24.44^\circ$, is the angle between the earth's axis of rotation and the perpendicular to its orbital plane. The other reason is the \emph{eccentricity} of the earth's orbit; the earth orbits the sun in an ellipse, not a perfect circle, and the eccentricity (about $0.0167$) measures the shape of the orbit (an eccentricity of $0$ corresponds to a perfect circle). But it's the obliquity of the earth's orbit that has the greater effect on the variation in the length of the solar day.

This activity is about understanding the solar day for a plane that orbits the sun in a perfect circle (the eccentricity of its orbit is zero) with an orbit that is perpendicular to its axis of rotation (the obliquity is zero). It is also about parameterizing the motion of a point at rest on the equator of that planet.



\begin{example} \label{Ec56r4322}
A planet with radius $b$ km revolves about its sun in a circular orbit of radius $a$ km at a constant speed. The equator of the planet lies in its orbital plane. There are four (solar) days in the planet's year and the planet rotates about its axis in the same sense as it revolves about its sun as illustrated below (where we see a cross-section of the sun and planet in the orbital plane).
\end{example}


\pdfOnly{Access Desmos interactive at \href{https://www.desmos.com/calculator/golluueo9z}{https://www.desmos.com/calculator/golluueo9z}}

 
\begin{onlineOnly}
    \begin{center}
\desmos{golluueo9z}{900}{600}
\end{center}
\end{onlineOnly}

\href{https://www.desmos.com/calculator/golluueo9z}{Rotation Rates of Planets}

Suppose the point $Q$, at rest relative to the planet's equator, has the position shown at the start of the calendar year (time $t=0$), when it is closest to the sun.

\begin{enumerate}
\item The first step is to count the number of rotations the planet makes about its axis during one year (the time it takes the planet  to make one revolution about its sun). To do this play or drag the slider $t_1$ (measured in years) in Line 2. The upper animation shows the planet rotating about its sun in the reference frame of the $xy$-coordinate system (fixed relative to the sun); the lower animation shows the motion of the point $Q$ relative to the $uv$-coordinate system with its origin at the planet's center). 

During a year with four (solar) days, the planet makes $\answer{5}$ rotations about its axis.

\item Now express the components of the vector
\[
  \overrightarrow{OP}  = \langle x , y\rangle =   \langle f_1(t) , g_1(t) \rangle
\]
giving the position (in kilometers) of the planet's center relative to the center of the sun in terms of time, measured in years (or fractions thereof) since the start of the year.

Since the orbital radius is $a$ km and the planet revolves about the sun at the rate of 
\[
    \omega = \answer{2\pi} \text{ radians/yr},
\]
the position vector is
\[
  \overrightarrow{OP}  = \langle x , y\rangle =  \langle \answer{a \cos (2\pi t)} , b \answer{\sin (2\pi t)} \rangle. 
\]

\item  Use the result of part (a) to express the components of the vector
 \[
  \overrightarrow{OQ} =  \langle u, v \rangle =   \langle f_2(t) , g_2(t) \rangle
\]
giving the position of point $Q$ relative to $O$ in terms of the number of years since the start of the year.

\begin{enumerate}
\item First express the position of $Q$ relative to $P$ in terms of the rotation angle $\alpha$ of the point $Q$ relative to $P$.

Keeping in mind the planet has radius $b$ km, the position vector is 
\begin{align*}
       \overrightarrow{PQ} &= \answer{\cos\alpha} \langle -b, 0 \rangle + \answer{\sin \alpha} \langle 0,-b  \rangle \\ 
                                     &=  \langle \answer{-b \cos \alpha}, \answer{-a \sin\alpha} \rangle . 
\end{align*}

\item Next express the rotation angle $\alpha$ in terms of $t$.

Since the planet rotates about its axis at the rate of $\answer{10\pi}$ radians/year,
\[
 \alpha = \answer{10\pi t}.
\]

\item Finally, express the position of $Q$ relative to $O$ in terms of $t$.

The position vector is
\begin{align*}
       \overrightarrow{OQ} &= \langle x , y \rangle \\
                                      &= \overrightarrow{OP} + \overrightarrow{OQ} \\
                                     &=  \langle \answer{a\cos t -b \cos (10\pi t)}, \answer{a \sin t - b\sin(10\pi t)} \rangle . 
\end{align*}


\end{enumerate}

\item Repeat parts (a)-(d) assuming the planet about its axis in the opposite sense that it revolves about its sun. (Drag the slider to $w=-1$ in Line 5 to see the motion.)
\end{enumerate}





\end{document}